% Advanced Mathematics
% Demonstrates: align, gather, theorem/proof, cases, split, and operator macros.
% Compile with: pdflatex main.tex

\documentclass[12pt]{article}
\usepackage[utf8]{inputenc}
\usepackage{amsmath}     % Core math environments
\usepackage{amssymb}     % Extra symbols (\mathbb, etc.)
\usepackage{amsthm}      % Theorem environments
\usepackage{mathtools}   % Extensions to amsmath (dcases, etc.)

% --- Theorem-like environments ---
\theoremstyle{plain}
\newtheorem{theorem}{Theorem}[section]
\newtheorem{lemma}[theorem]{Lemma}
\newtheorem{corollary}[theorem]{Corollary}

\theoremstyle{definition}
\newtheorem{definition}{Definition}[section]
\newtheorem{example}{Example}[section]

\theoremstyle{remark}
\newtheorem{remark}{Remark}

% --- Custom math operators ---
\DeclareMathOperator{\Tr}{Tr}       % Trace
\DeclareMathOperator{\diag}{diag}   % Diagonal

\title{Advanced Mathematics in \LaTeX{}}
\author{LaTeX Tutorial}
\date{\today}

\begin{document}
\maketitle
\tableofcontents
\newpage

% =============================================================================
\section{Multi-line Equations}
% =============================================================================

% --- align: each line independently aligned at & ---
\subsection{The \texttt{align} Environment}

Use \texttt{align} when each line should be numbered, and \texttt{align*}
to suppress numbering.

\begin{align}
    (a + b)^2 &= a^2 + 2ab + b^2       \label{eq:square-sum} \\
    (a - b)^2 &= a^2 - 2ab + b^2       \label{eq:square-diff} \\
    a^2 - b^2 &= (a+b)(a-b)            \label{eq:diff-squares}
\end{align}

From \eqref{eq:square-sum} and \eqref{eq:square-diff} we can derive
\eqref{eq:diff-squares}.

% --- gather: centered, no alignment column ---
\subsection{The \texttt{gather} Environment}

Use \texttt{gather} when equations should be centered but not aligned.

\begin{gather}
    \nabla \cdot \mathbf{E} = \frac{\rho}{\varepsilon_0} \\
    \nabla \cdot \mathbf{B} = 0 \\
    \nabla \times \mathbf{E} = -\frac{\partial \mathbf{B}}{\partial t} \\
    \nabla \times \mathbf{B} = \mu_0 \mathbf{J}
        + \mu_0 \varepsilon_0 \frac{\partial \mathbf{E}}{\partial t}
\end{gather}

% --- split inside equation: single number for multiple lines ---
\subsection{The \texttt{split} Environment}

Use \texttt{split} inside \texttt{equation} to get one number for a multi-line
derivation.

\begin{equation}
\begin{split}
    \int_0^\infty e^{-x^2}\,dx
        &= \frac{1}{2}\int_{-\infty}^{\infty} e^{-x^2}\,dx \\
        &= \frac{1}{2}\sqrt{\pi} \\
        &= \frac{\sqrt{\pi}}{2}
\end{split}
\label{eq:gaussian}
\end{equation}

% =============================================================================
\section{Cases and Piecewise Functions}
% =============================================================================

The \texttt{cases} environment (and \texttt{dcases} from \texttt{mathtools}
for display-style fractions):

\begin{equation}
    |x| =
    \begin{dcases}
        x  & \text{if } x \ge 0 \\
        -x & \text{if } x < 0
    \end{dcases}
\end{equation}

\begin{equation}
    f(n) =
    \begin{dcases}
        \frac{n}{2}   & \text{if } n \equiv 0 \pmod{2} \\[6pt]
        3n + 1         & \text{if } n \equiv 1 \pmod{2}
    \end{dcases}
\end{equation}

% =============================================================================
\section{Theorems, Lemmas, and Proofs}
% =============================================================================

\begin{definition}
A function $f: \mathbb{R} \to \mathbb{R}$ is \emph{continuous} at $a$ if
\[
    \lim_{x \to a} f(x) = f(a).
\]
\end{definition}

\begin{theorem}[Intermediate Value Theorem]
\label{thm:ivt}
If $f$ is continuous on $[a,b]$ and $k$ is between $f(a)$ and $f(b)$, then
there exists $c \in (a,b)$ such that $f(c) = k$.
\end{theorem}

\begin{proof}
Define $S = \{x \in [a,b] : f(x) \le k\}$. Since $S$ is bounded above and
non-empty, let $c = \sup S$. By continuity of $f$ at $c$, for every
$\varepsilon > 0$ there exists $\delta > 0$ such that $|x - c| < \delta$
implies $|f(x) - f(c)| < \varepsilon$. This forces $f(c) = k$.
\end{proof}

\begin{corollary}
Every polynomial of odd degree has at least one real root.
\end{corollary}

\begin{proof}
Follows from Theorem~\ref{thm:ivt} since
$\lim_{x \to \pm\infty} p(x) = \pm\infty$ for odd-degree $p$.
\end{proof}

\begin{lemma}[Triangle Inequality]
For all $x, y \in \mathbb{R}$:
\begin{equation}
    |x + y| \le |x| + |y|
\end{equation}
\end{lemma}

\begin{remark}
The proof of the IVT above is a sketch. A rigorous proof uses the
completeness of $\mathbb{R}$.
\end{remark}

% =============================================================================
\section{Matrices and Custom Operators}
% =============================================================================

A matrix with the custom $\Tr$ operator:
\begin{equation}
    A = \begin{pmatrix}
        a_{11} & a_{12} & a_{13} \\
        a_{21} & a_{22} & a_{23} \\
        a_{31} & a_{32} & a_{33}
    \end{pmatrix},
    \qquad
    \Tr(A) = \sum_{i=1}^{3} a_{ii}
\end{equation}

Diagonal matrix:
\begin{equation}
    D = \diag(\lambda_1, \lambda_2, \ldots, \lambda_n)
\end{equation}

\end{document}
